% ID: SEL-D0012-BD-E07
% ESTADO: cotejada. Segunda pasada visual completa y borrador compilado revisados.
% Fuente: Selectividad/matematicas_aplicadas/EXTRA_RESERVA_B_Modelo-6-Enunciados (Mat-Ap-2021).pdf
% SHA-256: 1752ec2f46b4ef2cc5d1c3346d69b79c2cc8a7aa12f6a0cd8f0d994b5bcda157
% Página: 2; bloque D; ejercicio original 7.

Para un estudio acerca del uso del transporte público en una ciudad,
se selecciona una muestra aleatoria de 500 individuos, obteniéndose que
175 de ellos lo usan.
\begin{enumerate}[label=\alph*)]
  \item \textbf{(1.5 puntos)} Halle un intervalo de confianza al 94\% para
  estimar la proporción real de individuos que usan el transporte público
  en esa ciudad.
  \item \textbf{(1 punto)} Manteniendo la proporción muestral,
  \textquestiondown{}cuántos individuos se deberían seleccionar como mínimo,
  para que, con un nivel de confianza del 97\%, la proporción muestral
  difiera de la proporción real a lo sumo en un 2\%?
\end{enumerate}
